\part{5G architecture and real-network integration}

\FigFiveG

\section{Understanding the 5G system}

\begin{qa}{575}{Explain the path from a user device to an external application through a 5G network.}
\begin{enumerate}
\item The UE registers through the gNB with the AMF (NAS over N1/N2; authentication via AUSF and UDM).
\item The UE requests a PDU session for an S-NSSAI and DNN; the SMF selects a UPF and installs packet rules over N4
(PFCP).
\item User packets go UE $\rightarrow$ (radio) $\rightarrow$ gNB $\rightarrow$ GTP-U tunnel on N3 $\rightarrow$ UPF
$\rightarrow$ N6 $\rightarrow$ data network.
\end{enumerate}
In our deployment the radio leg is simulated by UERANSIM.
\end{qa}

\begin{qa}{576}{What are the responsibilities of the UE, gNB, RAN and 5G Core?}
UE: the subscriber device and its NAS stack. gNB / RAN: radio access, radio scheduling, mobility, the N2/N3 interfaces.
5G Core: registration, sessions, policy, subscriber data, user-plane forwarding.
\end{qa}

\begin{qa}{577}{What are the roles of the AMF, SMF, UPF, PCF, UDM, AUSF and NRF?}
\begin{itemize}
\item AMF: access, mobility, NAS termination.
\item SMF: session management, UPF control (PFCP), IP allocation.
\item UPF: user-plane forwarding, GTP-U, QoS enforcement (QER), usage reporting.
\item PCF: policy (QoS and charging rules).
\item UDM: subscription data and credentials.
\item AUSF: authentication.
\item NRF: network-function discovery.
\end{itemize}
(Plus SCP, NSSF, BSF, UDR in our deployment.)
\end{qa}

\begin{qa}{578}{Control-plane signalling vs user-plane traffic?}
Control: NAS and NGAP (N1/N2), service-based HTTP/2 between network functions, PFCP (N4). User: GTP-U-encapsulated
packets (N3) and plain IP (N6). CARS-SF acts only on the user plane (UPF egress).
\end{qa}

\begin{qa}{579}{What is a PDU session?}
A logical connection between a UE and a data network (DNN) through a UPF, with an IP address and one or more QoS
flows, created for a given S-NSSAI.
\end{qa}

\begin{qa}{580}{What is a QoS flow, and how does it relate to a PDU session?}
The finest QoS granularity in 5G: packets are mapped by filters (SDF templates) to a QoS flow (identified by QFI)
within a PDU session, each with a 5QI and parameters (GBR/MBR, ARP).
\end{qa}

\begin{qa}{581}{What are 5QI values, and how do they relate to QoS treatment?}
Standardised indices (3GPP TS 23.501) mapping to a resource type (GBR, non-GBR, delay-critical), priority, packet delay
budget and error rate. For example 5QI 1 is conversational voice and 5QI 9 is best-effort default.
\honest{All our subscribers use 5QI 9. Differentiation comes from CARS-SF's UPF qdiscs, not 5QI.}
\end{qa}

\begin{qa}{582}{What is the difference between a network slice and a QoS flow?}
A slice is an end-to-end logical network (S-NSSAI) with its own network-function instances, policy or resources. A
QoS flow is a traffic treatment class inside a PDU session within a slice. CARS-SF's ``service slice'' maps to an
S-NSSAI plus DNN in the 5G deployment.
\end{qa}

\begin{qa}{583}{What is an S-NSSAI?}
Single Network Slice Selection Assistance Information: identifies a slice as SST plus an optional SD.
\end{qa}

\begin{qa}{584}{What are SST and SD?}
SST: Slice/Service Type (1 eMBB, 2 URLLC, 3 MIoT standardised). SD: Slice Differentiator, distinguishing slices of the
same SST. We use healthcare SST 2/SD 000001, education 1/000002, agriculture 3/000003, general 1/000004.
\end{qa}

\begin{qa}{585}{What is the difference between eMBB, URLLC and mMTC?}
Enhanced Mobile Broadband (high throughput), Ultra-Reliable Low-Latency Communication (strict latency and
reliability), massive Machine-Type Communication (many low-rate devices). These are technology classes; CARS-SF argues
for service classes on top (the 3GPP baseline in the main study maps applications to them).
\end{qa}

\begin{qa}{586}{What is the RAN's role in allocating radio resources?}
The gNB's MAC scheduler allocates PRBs per TTI among UEs and QoS flows (and slices, with RAN slicing). This is where
most capacity contention lives in real cellular networks. It is absent from our testbed.
\end{qa}

\begin{qa}{587}{Why can transport-layer bandwidth allocation not, by itself, guarantee radio-layer QoS?}
If the radio is the bottleneck (poor coverage, cell congestion), packets queue or are lost in the RAN before reaching
the backhaul. Shaping at the UPF or rural uplink then shapes a link that is not the constraint. CARS-SF's allocation is
meaningful when the backhaul or transport is the bottleneck, which is our stated rural assumption.
\end{qa}

\begin{qa}{588}{What is the role of edge computing in rural connectivity?}
Hosting critical services (clinic, agri) at the edge keeps them reachable and fast when the backhaul fails or degrades.
Our topology shows this: edge-hosted healthcare survives a core switch failure, while core-hosted services cannot.
\end{qa}

\begin{qa}{589}{How could backhaul constraints interact with radio-resource constraints?}
They form two bottlenecks in series. If they are allocated inconsistently (the RAN favours education, the backhaul
favours healthcare), the tighter one dominates and capacity is wasted. Joint slicing (RAN plus transport) would need a
shared policy, e.g.\ the same weights pushed to RAN-slice quotas and to transport HTB.
\end{qa}

\begin{qa}{590}{Which of these concepts are implemented in CARS-SF, and which are background?}
\textbf{Implemented and run} (L3, simulated RAN):
\begin{itemize}
\item registration, PDU sessions;
\item four S-NSSAIs and DNNs with their own address pools;
\item GTP-U user plane through the UPF;
\item CARS-SF enforcement on UPF N6 and \code{ogstun}.
\end{itemize}
\textbf{Background only:} 5QI differentiation, PCF policy, RAN scheduling and slicing, QoS flows beyond the default,
mobility.
\end{qa}

\section{Open5GS, UERANSIM and deployment claims}

\begin{qa}{591}{What is Open5GS?}
An open-source implementation of the 4G EPC and 5G Core (AMF, SMF, UPF, NRF, SCP, AUSF, UDM, UDR, PCF, NSSF, BSF),
used here as version 2.7.2 in a single container.
\end{qa}

\begin{qa}{592}{What is UERANSIM?}
An open-source 5G UE and gNB simulator: it implements NAS, RRC-level procedures and GTP-U, but not a physical radio
layer. Version 3.2.6.
\end{qa}

\begin{qa}{593}{Which components does each provide?}
Open5GS: the full 5GC network functions. UERANSIM: \code{nr-gnb} (NGAP over SCTP, GTP-U) and \code{nr-ue}
(registration, PDU sessions, TUN interface \code{uesimtun0}).
\end{qa}

\begin{qa}{594}{What is needed for a working UE-to-core-to-data-network path?}
\begin{itemize}
\item Matching PLMN (999/70), TAC and slices in gNB, AMF and NSSF.
\item Subscribers in MongoDB (keys, OPc, slice and DNN).
\item SMF session pools per DNN; UPF with \code{ogstun} carrying each slice gateway.
\item The UPF running as root (TUN creation).
\item SCTP to the AMF (start order mattered: the gNB does not retry).
\item Routes from UEs via \code{uesimtun0}, and from the data network back to the slice pools via the UPF's N6.
\end{itemize}
All of it is in \code{deployment/5g/} and was brought up and verified (ping in both directions).
\end{qa}

\begin{qa}{595}{Where would CARS-SF sit in that architecture?}
As an application function at the edge. Today it is a controller that enforces on the UPF's N6 and \code{ogstun}
egress via tc. A standards-aligned version would sit behind the PCF or NEF (an AF influencing policy) and steer QERs.
\end{qa}

\begin{qa}{596}{Which actual interfaces would CARS-SF use to observe and influence the 5G system?}
\textbf{Now:} Linux interfaces on the UPF host (tc stats, counters) and captures on N6 and the UE TUNs.
\textbf{Standards path:}
\begin{itemize}
\item NWDAF or UPF usage reports (observation);
\item Npcf / N5 via the AF, and NEF for policy;
\item SMF $\rightarrow$ UPF QERs (MBR/GBR) over N4;
\item RAN slice quotas via O-RAN A1/E2 for radio.
\end{itemize}
\end{qa}

\begin{qa}{597}{How would an allocation decision become a supported 5G configuration?}
Map slice allocations to session or QoS-flow AMBR/MBR/GBR and QER updates (through the PCF, as policy for each S-NSSAI
and DNN), and to RAN-slice PRB quotas. Our tc shaping on the UPF host is a pragmatic stand-in with the same effect
on the N6 bottleneck.
\end{qa}

\begin{qa}{598}{Which 5G functions would need modification?}
None for the demonstrated mode (tc on the UPF host is outside 5G functions). For a standards-aligned integration: none
in principle, using the PCF with AF influence and QER-based enforcement in the UPF. In practice, Open5GS's support for
dynamic QER updates driven by an AF would need checking or extending, and the RAN side needs an E2-capable gNB (not
UERANSIM).
\end{qa}

\begin{warnbox}{Reconstructed questions}
The questionnaire was truncated at Q598. Questions 598a onward cover the rest of the 5G section, experiments and
critical evaluation, written to match the questionnaire's intent and its stated headings (``Experiments \textperiodcentered{}
Critical evaluation'').
\end{warnbox}

\begin{qa}{598a}{What exactly did you run on the 5G core, and at which evidence level?}
\begin{itemize}
\item Open5GS with four slice UEs (UERANSIM) and PDU sessions on their own S-NSSAI/DNN pools.
\item CARS-SF traffic generators in sidecars sharing each UE's namespace.
\item A 10\,Mb/s UPF bottleneck shaped by CARS-SF, compared with no slicing, in the disaster and school contexts.
\item Captures on \code{uesimtun0} (raw IP) and on the data network, matched with the same code.
\end{itemize}
\Lthreefg. It is not L4: the radio is simulated.
\end{qa}

\begin{qa}{598b}{What did the 5G experiment show?}
Disaster context, 10\,Mb/s UPF:
\begin{itemize}
\item healthcare p95 \NFgDisasterCarsHcP\,ms with CARS-SF vs \NFgDisasterNoneHcP\,ms without;
\item healthcare SLA-met share \NFgDisasterCarsHcSla{} vs \NFgDisasterNoneHcSla;
\item agriculture p95 \NFgDisasterCarsAgP{} vs \NFgDisasterNoneAgP\,ms.
\end{itemize}
Trade-off: weighted throughput satisfaction \NFgDisasterCarsUtil{} vs \NFgDisasterNoneUtil, because fq\_codel drops
some of healthcare's bursty video in exchange for low delay.
School: healthcare p95 \NFgSchoolCarsHcP{} vs \NFgSchoolNoneHcP\,ms. Healthcare is uncontended there; the report marks
H8 (``lower in every scenario'') \textbf{supported}, but this difference is about 1\,ms with $n=3$ and had the opposite
sign before the audit, so I claim only the disaster effect; weighted utility still improves
(\NFgSchoolCarsUtil{} vs \NFgSchoolNoneUtil).
\end{qa}

\begin{qa}{598c}{Why a 10\,Mb/s bottleneck on the 5G core, not 100\,Mb/s?}
At 20\,Mb/s, healthcare showed about 390\,ms p95 in \emph{both} modes. The queue was in UERANSIM's user-space GTP path,
not at the UPF. At 10\,Mb/s the UPF is the contention point, which is what the experiment needs. This is a documented
limitation of UERANSIM, not of CARS-SF.
\end{qa}

\begin{qa}{598d}{Does slicing at the UPF with tc constitute 5G network slicing?}
Partially. Slice \emph{identity} is genuine 5G (S-NSSAI, DNN, PDU sessions through the 5GC); slice \emph{resource
enforcement} is CARS-SF's tc on the UPF host, not 3GPP QoS (5QI, QER, PCF). Say exactly that.
\end{qa}

\begin{qa}{598e}{How would you reach L4 evidence?}
Replace UERANSIM with an SDR-based gNB (e.g.\ a USRP B210 with an open-source RAN stack) and commercial UEs on a
private PLMN, as in the repository's existing 5G testbed proposal (\code{existing/5G_TestBed}). Then run the
fiber-cut and disaster scenarios over the air, with a measured backhaul bottleneck.
\end{qa}

\begin{qa}{598f}{What were the deployment pitfalls, and why do they matter for the defense?}
\begin{itemize}
\item The image runs as non-root, so TUN creation failed.
\item A read-only \code{/proc/sys}.
\item MongoDB readiness.
\item YAML indentation of \code{s_nssai}.
\item The gNB connected 16\,ms before the AMF listened and did not retry.
\item Captures on a TUN device use LINKTYPE\_RAW (101), which our parser initially dropped.
\end{itemize}
Each was found and fixed by running end to end. That is evidence the result is from a working system, not a diagram.
\end{qa}

\begin{qa}{598g}{What is the evidence level of the OVS QoS backend and of real-OVS operation generally?}
Real OVS 3.3.9 with a genuine \code{OVSBridge}, run in a privileged container (\Lthreeovs):
\begin{itemize}
\item a fiber-cut run: recovery \NOvsMttrAfter\,s, outage \NOvsOutageAfter\,s;
\item a school run with the \code{linux-htb} backend: no false failover after the congestion fix, validation 93.
\end{itemize}
\honest{Only a handful of runs. The 340-run matrix used Linux bridges.}
\end{qa}
